\documentclass[10pt,a4paper]{article}

\usepackage[a4paper,margin=0.7in]{geometry}
\usepackage{booktabs}
\usepackage{amsmath,amssymb}
\usepackage{array}
\usepackage{longtable}
\usepackage{graphicx}

\title{OR114-2 Midterm Project}
\author{}
\date{}


\begin{document}
\maketitle

\section*{Problem 1}
\subsection*{1. Mathematical Formulation}

\textbf{Sets}

\begin{itemize}
    \item $C$ : Set of vehicles.
    \item $O$ : Set of customer orders.
\end{itemize}

\begin{itemize}
    \item $p_i$ : Pickup station of order $i$.
    \item $r_i$ : Return station of order $i$.
    \item $s_i$ : Pick-up time of order $i$.
    \item $e_i$ : Return time of order $i$.
    \item $l_i$ : Requested vehicle level for order $i$.
    \item $R_i$ : Revenue generated by order $i$.
\end{itemize}

\begin{itemize}
    \item $h_c$ : Initial station of vehicle $c$.
    \item $q_c$ : Vehicle level of vehicle $c$.
\end{itemize}

\textbf{Parameters}

\begin{itemize}
    \item $A^0$: set of feasible initial assignments
    \item $A$: set of feasible order transitions
    \item $m_{ci}$: relocation time from vehicle c to order i 
    \item $m_{ij}$: relocation time between orders i and j 
    \item $B$: relocation time budget
\end{itemize}

The feasible sets $A^0$ and $A$ are defined by IEDO timing constraints, where a transition $(i, j)$ is valid only if:
\[ RE_i + 240 + T_{ab} \le PU_j - 30 \]
\hspace{1em} where $RE_i$, $T_{ab}$, and $PU_j$ represent return time, travel time, and pickup time, respectively.

\textbf{Decision Variables}

\[
y_i =
\begin{cases}
1, & \text{if order } i \text{ is accepted}\\
0, & \text{otherwise}
\end{cases}
\]

\[
x_{ci} =
\begin{cases}
1, & \text{if vehicle } c \text{ starts with order } i\\
0, & \text{otherwise}
\end{cases}
\]

\[
z_{cij} =
\begin{cases}
1, & \text{if vehicle } c \text{ serves order } j \text{ after } i\\
0, & \text{otherwise}
\end{cases}
\]

\vspace{0.2cm}

\textbf{Objective Function}

\[
\max \sum_{i \in O} R_i y_i
\]

\vspace{0.2cm}

\textbf{Constraints}

\textbf{(1) Order assignment constraint}

\[
\sum_{(c,i)\in A^0} x_{ci}
+
\sum_{(c,j,i)\in A} z_{cji}
=
y_i
\qquad \forall i \in O
\]

\textbf{(2) Vehicle starting constraint}

\[
\sum_{i:(c,i)\in A^0} x_{ci}
\le 1
\qquad \forall c \in C
\]

\textbf{(3) Flow conservation constraint}

\[
\sum_{j:(c,i,j)\in A} z_{cij}
\le
x_{ci}
+
\sum_{k:(c,k,i)\in A} z_{cki}
\qquad \forall c \in C, i \in O
\]

\textbf{(4) Relocation budget constraint}

\[
\sum_{(c,i)\in A^0} m_{ci}x_{ci}
+
\sum_{(c,i,j)\in A} m_{ij}z_{cij}
\le B
\]

\vspace{0.3cm}

\subsection*{2. Optimal Solution for Instance 05}

\hspace{2em} The optimization model accepts 8 out of 10 orders, including Orders 3, 4, 5, 6, 7, 8, 9, and 10. Orders 1 and 2 are rejected. The accepted sales revenue is 117,000, and the final profit after rejection compensation costs is 106,800. The total relocation time used is 1,080 minutes, which satisfies the relocation budget limit ($1080 \le 1200$). Only one vehicle upgrade is required in the final solution, where a level-2 vehicle is assigned to Order 6, which requests a level-1 vehicle.

\subsection*{3. Operational Plan}

\small

\begin{longtable}{cccccc}
\toprule
Order & Car & Upgrade & Pickup $\rightarrow$ Return & Rental Period & Revenue \\
\midrule
3 & 3 & No & 5 $\rightarrow$ 1 & 01/02 03:30 -- 01/04 11:30 & 5,600 \\
4 & 1 & No & 1 $\rightarrow$ 2 & 01/03 12:00 -- 01/05 15:00 & 5,100 \\
5 & 6 & No & 6 $\rightarrow$ 6 & 01/01 00:00 -- 01/04 20:00 & 36,800 \\
6 & 4 & Yes & 3 $\rightarrow$ 4 & 01/04 23:00 -- 01/05 14:00 & 1,500 \\
7 & 5 & No & 2 $\rightarrow$ 1 & 01/02 03:30 -- 01/04 23:30 & 27,200 \\
8 & 6 & No & 6 $\rightarrow$ 6 & 01/05 04:30 -- 01/05 14:30 & 4,000 \\
9 & 4 & No & 3 $\rightarrow$ 4 & 01/01 19:30 -- 01/04 15:30 & 27,200 \\
10 & 2 & No & 4 $\rightarrow$ 2 & 01/01 06:00 -- 01/05 06:00 & 9,600 \\
\bottomrule
\end{longtable}

\vspace{0.2cm}

\subsection*{4. Vehicle Relocation Plan}

\begin{table}[h]
\centering
\small
\begin{tabular}{cccccc}
\toprule
Car & From & To & Depart & Arrive & Minutes \\
\midrule
2 & 2 & 4 & 01/01 00:00 & 01/01 04:00 & 240 \\
3 & 3 & 5 & 01/01 00:00 & 01/01 03:30 & 210 \\
4 & 5 & 3 & 01/01 00:00 & 01/01 03:30 & 210 \\
4 & 4 & 3 & 01/04 19:30 & 01/04 22:30 & 180 \\
5 & 4 & 2 & 01/01 00:00 & 01/01 04:00 & 240 \\
\bottomrule
\end{tabular}
\end{table}

\section*{Problem 2}

For the large hidden instances, an exact IP model may be too slow.  We therefore
implemented a heuristic algorithm that always returns a feasible assignment and
relocation plan within the three-minute grading limit.  The submitted algorithm
has two main stages.  First, it quickly constructs two feasible plans using two
different heuristics, called Algo1 and Algo5 in our implementation.  It keeps the
plan with the larger profit as the initial solution.  Second, it applies a
batch-based local integer programming improvement procedure.  In each batch, a
small number of cars are temporarily released, together with the orders currently
served by those cars, and a small IP is solved only for this local subproblem.
The new local plan is accepted only if it improves the total profit and still
satisfies the moving-time budget.

Algo1 is a greedy route-insertion method.  Orders are sorted by high revenue
first, then by earlier pickup time.  For each order, the algorithm tries every
compatible car, including the free upgrade rule that a level-$l$ order may be
served by a level-$l$ or level-$(l+1)$ car.  The order is inserted into the
chronological route of the car if both adjacent transitions are feasible.  Among
all feasible insertions, Algo1 chooses the one with the smallest additional
moving time, then the smallest upgrade penalty and idle time.

Algo5 is a demand-aware look-ahead heuristic.  It divides the planning horizon
into six-hour buckets and builds a future demand score for each station and car
level.  When assigning a car to an order, Algo5 scores the decision by combining
the immediate profit benefit, relocation time, idle time, upgrade penalty, the
future value of the return station, and the opportunity cost of removing the car
from its current station.  After this greedy pass, it performs several bounded
repairs, such as shortage-bucket repair, route insertion, and one-removal
replacement, to rescue valuable rejected orders without violating feasibility.

The final algorithm, implemented in \texttt{algo\_union.py}, uses the following
rule:
\[
\texttt{pre\_build}=\arg\max\{\text{profit(Algo1)},\text{profit(Algo5)}\}.
\]
Because both initial heuristics are fast and feasible, taking their better plan
is a low-risk way to combine two complementary decision rules.  The resulting
plan is then improved by \texttt{small\_ip\_improve}, which reuses the local IP
repair idea from Algo4.  A batch of cars is selected according to route
efficiency.  The orders on these cars are released; high-revenue compatible
orders are added as candidates; then a local network-flow IP is solved over this
small set.  The IP uses the same type of start variables, link variables, and
acceptance variables as in Problem 1, but only for the selected cars and
candidate orders.  The local repair is controlled by a wall-clock time limit, so
the algorithm can safely stop and return the best feasible plan found so far.

\section*{Problem 3}

The intuition behind our heuristic is that the global problem is hard mainly
because it couples many vehicle routes through the shared relocation budget.  A
pure greedy method is fast, but it may spend scarce relocation minutes on early
orders and block better future opportunities.  A full IP sees these tradeoffs
but is too large.  Our algorithm therefore combines a fast global construction
with repeated small exact repairs.

\begin{enumerate}
    \item \textbf{Build two initial plans.}  Algo1 emphasizes high immediate
    revenue and low incremental relocation cost.  Algo5 emphasizes future
    demand balance by using station-level look-ahead scores.  These two
    heuristics tend to make different mistakes, so we keep the more profitable
    of the two.
    \item \textbf{Convert the assignment into car routes.}  This allows us to
    calculate route-level moving time and profit, and to remove or repair only a
    few routes at a time.
    \item \textbf{Release a small batch of cars.}  Cars are sampled with higher
    probability when their routes have lower efficiency.  This focuses repair
    effort on the weakest part of the current solution.
    \item \textbf{Solve a local IP.}  For the selected cars, we include their
    released orders and a capped list of compatible high-value candidate orders.
    The local IP chooses a feasible set of routes for only this batch under the
    remaining relocation budget.
    \item \textbf{Accept only improving repairs.}  If the repaired routes
    increase profit and keep total moving time within $B$, they replace the old
    routes.  Otherwise the old feasible solution is kept.
\end{enumerate}

The procedure is summarized below.

\[
\begin{array}{ll}
\textbf{Input:} & \text{one instance file, time limit }T.\\
1. & (A_1,R_1)\leftarrow\text{Algo1(instance)}.\\
2. & (A_5,R_5)\leftarrow\text{Algo5(instance)}.\\
3. & A\leftarrow \text{the assignment among }A_1,A_5\text{ with larger profit}.\\
4. & \textbf{while time remains do}\\
5. & \quad \text{sample a small batch of cars from current routes};\\
6. & \quad \text{release their served orders and build candidate order set};\\
7. & \quad \text{solve the local IP with a short per-IP time limit};\\
8. & \quad \text{if the local solution improves profit, update the routes};\\
9. & \textbf{return } A \text{ and the corresponding relocation list.}
\end{array}
\]

Let $n_C$ be the number of cars, $n_K$ the number of orders, and $n_S$ the
number of stations.  The construction heuristics examine compatible car-order
pairs and route insertions, so their dominant cost is roughly
$O(n_C n_K+n_K r)$, where $r$ is the number of route positions checked in the
repair rules.  In the local IP stage, suppose each batch contains $b$ cars and
at most $m$ candidate orders.  The local network has $O(bm)$ start arcs and
$O(bm^2)$ order-to-order arcs in the worst case, so the local IP has
$O(bm^2)$ binary variables and constraints.  In practice, $b$ and $m$ are capped
by parameters and the entire loop is bounded by a wall-clock time limit.  Thus
the algorithm scales to large instances by spending exact optimization effort
only on small neighborhoods of the current solution.

\section*{Problem 4}

Following the project instruction, we designed a numerical study with explicit
factors and scenarios rather than relying only on the five public instances.  The
random instances are generated by \texttt{data/generate\_code.py}.  The planning
horizon always starts at 2023/01/01 00:00.  For each instance, the number of
stations is fixed at $n_S=100$, the number of levels is fixed at $n_L=10$, and
the number of planning days $n_D$ is sampled uniformly from 7 to 100.  Hourly
rates are generated as a strictly increasing sequence over levels.  Moving time
is symmetric, with $T_{ii}=0$ and $T_{ij}=90$ minutes for $i\ne j$ in the
generated study instances.

\begin{table}[h]
\centering
\small
\begin{tabular}{clcccc}
\toprule
Scenario & Main factor & Ratio $n_K:n_C$ & Levels & Stations & Time/Budget \\
\midrule
S1 & Baseline & 1:1 & random & random & random, $B=10^6$\\
S2 & Ultra-low load & 1:10 & random & random & random, $B=10^6$\\
S3 & Low load & 1:3 & random & random & random, $B=10^6$\\
S4 & High load & 3:1 & random & random & random, $B=10^6$\\
S5 & Level mismatch & 1:1 & 8:2 skew & random & random, $B=10^6$\\
S6 & Forced relocation & 1:1 & random & odd/even & random, $B=300n_D$\\
S7 & One hub & 1:1 & random & hub 1G & random, $B=300n_D$\\
S8 & Two hubs & 1:1 & random & hub 2G & random, $B=300n_D$\\
S9 & Weekend effect & 1:1 & random & random & weekend, $B=10^6$\\
S10 & One peak & 1:1 & random & random & one peak, $B=10^6$\\
S11 & Three peaks & 1:1 & random & random & three peaks, $B=10^6$\\
S12 & No relocation budget & 1:1 & random & random & random, $B=0$\\
S13 & Tight budget & 1:1 & random & random & random, $B=30n_D$\\
S14 & Moderate budget & 1:1 & random & random & random, $B=300n_D$\\
\bottomrule
\end{tabular}
\end{table}

For the final detailed experiment, we selected five representative hard
scenarios: S6, S8, S10, S13, and S14.  These cover spatial imbalance, hub
returns, temporal bursts, and relocation-budget sensitivity.  For each selected
scenario we generated 128 random instances.  For each instance, we solved the IP
model from Problem 1 with Gurobi to obtain an optimal benchmark, then compared
our \texttt{algo\_union} heuristic with a simple baseline, \texttt{algo\_naive}.
The naive baseline processes orders chronologically and assigns each order to
the feasible car requiring the least relocation time; if no such car exists, the
order is rejected.  The optimality gap is reported as
\[
\text{gap}=\frac{\text{IP profit}-\text{algorithm profit}}{|\text{IP profit}|}.
\]
Lower is better, and zero means that the heuristic matched the IP profit.

\begin{table}[h]
\centering
\small
\begin{tabular}{lrrrr}
\toprule
Scenario & Naive mean gap & Naive std. & Union mean gap & Union std.\\
\midrule
S6  & 59.88\% & 222.65\% & 0.00\% & 0.00\%\\
S8  & 97.53\% & 550.77\% & 0.68\% & 5.36\%\\
S10 & 27.60\% & 12.85\% & 2.16\% & 2.01\%\\
S13 & 1063.61\% & 3967.26\% & 29.07\% & 140.57\%\\
S14 & 60.83\% & 233.16\% & 1.10\% & 4.66\%\\
\bottomrule
\end{tabular}
\end{table}

The histogram figures are generated under the scenario-study result directory.
It contains one combined SVG histogram and one SVG histogram for each selected
scenario.  The full raw experiment, including the 640 generated instances,
optimal plans, run-time logs, per-instance benchmark rows, and histograms, is
stored in the compressed artifact \texttt{scenario\_study\_128\_results.tgz}.

The results show that \texttt{algo\_union} is consistently much closer to the IP
benchmark than the naive chronological rule.  On S6, the union heuristic matched
the IP solution on all 128 instances.  On S8, S10, and S14, the average gap
remained below 2.2\%, while the naive baseline had gaps between 27.6\% and
97.5\%.  S13 is the hardest case because the relocation budget is very tight and
the IP profit is often negative or close to zero, which makes percentage gaps
large and volatile.  Even in this case, \texttt{algo\_union} reduced the average
gap from 1063.61\% for the naive baseline to 29.07\%.  This supports the main
design idea: a fast greedy construction is useful, but local batch IP repair is
crucial when relocation budget or spatial imbalance makes early greedy choices
expensive.

\end{document}
